\documentclass[12pt,a4paper]{article}

\usepackage[a4paper, margin=1in]{geometry}
\usepackage{titlesec}
\usepackage{graphicx}
\usepackage{float}
\usepackage{longtable}
\usepackage{array}
\usepackage{setspace}
\usepackage{hyperref}
\usepackage{enumitem}
\usepackage[table]{xcolor}
\usepackage{fancyhdr}
\usepackage{tocloft}
\usepackage{booktabs}
\usepackage{tabularx}

% ---------- Color Theme ----------
\definecolor{Primary}{HTML}{0B3C5D}
\definecolor{Secondary}{HTML}{1D70A2}
\definecolor{Accent}{HTML}{F39C12}
\definecolor{SoftBlue}{HTML}{EAF4FB}
\definecolor{SoftGray}{HTML}{F5F7FA}
\definecolor{SoftGreen}{HTML}{E9F7EF}

\hypersetup{
    colorlinks=true,
    linkcolor=Secondary,
    urlcolor=Secondary,
    pdftitle={Jobra Hospital Management System Report},
    pdfauthor={CSE Project Team},
    pdfsubject={DBMS Full Stack Project Report}
}

\pagestyle{fancy}
\fancyhf{}
\lhead{\textcolor{Primary}{\textbf{Project Report}}}
\rhead{\textcolor{Primary}{\textbf{Jobra Hospital Management System}}}
\cfoot{\textcolor{Secondary}{\thepage}}
\renewcommand{\headrulewidth}{0.8pt}
\renewcommand{\headrule}{\hbox to\headwidth{\color{Secondary}\leaders\hrule height \headrulewidth\hfill}}

\setstretch{1.25}
\setlength{\parskip}{0.35em}
\setlength{\parindent}{0pt}

\titleformat{\section}{\Large\bfseries\color{Primary}}{\thesection.}{0.5em}{}
\titleformat{\subsection}{\large\bfseries\color{Secondary}}{\thesubsection}{0.5em}{}
\titleformat{\subsubsection}{\normalsize\bfseries\color{Primary}}{\thesubsubsection}{0.5em}{}

\newcommand{\infobox}[1]{
    \vspace{0.4em}
    \noindent\colorbox{SoftBlue}{\parbox{\dimexpr\linewidth-2\fboxsep\relax}{#1}}
    \vspace{0.4em}
}

\newcommand{\placeholder}[2]{
    \begin{figure}[H]
        \centering
        \fcolorbox{Secondary}{SoftGray}{\parbox[c][5cm][c]{0.9\linewidth}{\centering\textcolor{Primary}{\textbf{#1}}\\[0.4em]\textcolor{Secondary}{#2}}}
        \caption{#1}
    \end{figure}
}

\begin{document}

% --------------------- COVER PAGE ---------------------
\begin{titlepage}
    \centering
    \vspace*{1.2cm}
    {\Huge\color{Primary}\textbf{Final Deliverable: Part-2}}\\[0.25cm]
    {\Large\color{Secondary}\textbf{Project Report and Deployment Documentation}}\\[1.4cm]
    
    {\Huge\textbf{\textcolor{Primary}{Jobra Hospital Management System}}}\\[0.6cm]
    {\large\textcolor{Secondary}{A Multi-Role Web-Based Hospital Service Platform}}\\[1.2cm]
    
    \rowcolors{2}{SoftGray}{white}
    \begin{tabular}{>{\bfseries}p{5.5cm} p{6.5cm}}
        Team Member Name & Tanjim Tajwar Arnab \\
        Student ID & 22701066 \\
        Team Member Name & Sabrina Sultana \\
        Student ID & 22701067 \\
        Team Member Name & Hafiz Hasnat Sifar Jami \\
        Student ID & 22701068 \\
        Team Member Name & Muznabin Ahmed \\
        Student ID & 22701069 \\
    \end{tabular}
    
    \vfill
    
    \infobox{
    \textbf{Course:} Database Management Systems / Web Engineering \hfill
    \textbf{Semester:} 6th\\
    \textbf{Department:} Computer Science and Engineering\\
    \textbf{Project Type:} Full Stack DBMS Application\\
    \textbf{Submission Date:} \today
    }
    
    \vfill
    {\small\textcolor{Secondary}{Prepared for academic evaluation and demonstration.}}
\end{titlepage}

% --------------------- PRELIMINARY PAGES ---------------------
\tableofcontents
\newpage

\section*{Executive Abstract}
\addcontentsline{toc}{section}{Executive Abstract}
The \textbf{Jobra Hospital Management System} is a role-based, full-stack web application developed to digitize common operational workflows of a hospital or clinic. The project combines a React frontend, a Node.js/Express backend, and a MySQL relational database to provide a complete data-driven platform for \textbf{Admin}, \textbf{Doctor}, and \textbf{Patient} users.

The system supports the complete lifecycle of healthcare interaction in this project scope: user registration and login, doctor-patient appointment management, clinical report writing, complaint handling, and profile management for each role. The solution emphasizes database normalization, role-based access control, and practical API integration.

\infobox{
\textbf{Core Outcome:} The project successfully demonstrates real-world DBMS concepts through multi-table relational data, role-aware business logic, CRUD operations, and secure authentication flow.
}

\newpage

% --------------------- 1. INTRODUCTION ---------------------
\section{Introduction}
\subsection{Background}
Manual hospital operations often lead to duplicated data entry, delayed communication, and inefficient record tracking. A digital system is required to centralize operational data and make patient services faster and more reliable.

\subsection{Problem Statement}
Hospitals need a secure and role-specific system where:
\begin{itemize}[leftmargin=1.2cm]
    \item Patients can access services online without physical queue dependency.
    \item Doctors can manage appointments and medical notes efficiently.
    \item Admin users can monitor and control platform-wide operations.
\end{itemize}

\subsection{Proposed Solution}
This project provides a web-based solution with:
\begin{itemize}[leftmargin=1.2cm]
    \item Role-based dashboards and protected routes.
    \item Centralized MySQL data storage.
    \item REST APIs for modular data operations.
    \item Secure login and token-based session flow.
\end{itemize}

% --------------------- 2. PROJECT GOALS ---------------------
\section{Project Goals and Objectives}
\begin{itemize}[leftmargin=1.2cm]
    \item Build a centralized hospital service management platform.
    \item Implement role-based access for Admin, Doctor, and Patient users.
    \item Allow patients to book, track, and cancel appointments.
    \item Enable doctors to update appointment status and write medical reports.
    \item Enable patients to submit complaints and feedback.
    \item Provide admin-level visibility across users, appointments, and complaints.
    \item Demonstrate practical full-stack and DBMS integration in an academic project.
\end{itemize}

% --------------------- 3. TECH STACK ---------------------
\section{Technology Stack}
\subsection{Frontend Layer}
\begin{itemize}[leftmargin=1.2cm]
    \item React (component-based SPA)
    \item React Router (client-side routing and route protection)
    \item Axios (API communication)
    \item Context API (authentication state handling)
    \item Custom CSS by module/role
\end{itemize}

\subsection{Backend Layer}
\begin{itemize}[leftmargin=1.2cm]
    \item Node.js runtime
    \item Express.js framework
    \item JWT for token-based authentication
    \item bcrypt for password hashing
    \item CORS and dotenv configuration
\end{itemize}

\subsection{Data Layer}
\begin{itemize}[leftmargin=1.2cm]
    \item MySQL (relational database)
    \item mysql2 package for query execution
    \item SQL schema file for reproducible database setup
\end{itemize}

% --------------------- 4. SYSTEM ARCHITECTURE ---------------------
\section{System Architecture}
\subsection{Architectural Pattern}
The implementation follows a layered full-stack architecture:
\begin{enumerate}[leftmargin=1.2cm]
    \item \textbf{Frontend UI Layer} (React pages/components).
    \item \textbf{Service Layer} (Axios services for API abstraction).
    \item \textbf{Backend Route Layer} (Express route definitions).
    \item \textbf{Controller/Model Layer} (business logic + SQL access).
    \item \textbf{MySQL Data Layer} (persistent relational storage).
\end{enumerate}

\subsection{Access Control Design}
\begin{itemize}[leftmargin=1.2cm]
    \item Middleware verifies token authenticity.
    \item Role-check middleware enforces endpoint-level permission.
    \item Frontend protected route component prevents unauthorized UI access.
\end{itemize}

\infobox{
\textbf{Architecture Advantage:} Separation of concerns keeps features modular, easier to maintain, and suitable for incremental extension in future semesters.
}

% --------------------- 5. USER ROLES ---------------------
\section{User Roles and Responsibilities}
\rowcolors{2}{SoftGray}{white}
\begin{tabularx}{\linewidth}{>{\bfseries}p{2.5cm} X}
\toprule
Role & Responsibilities \\
\midrule
Admin & View and manage system users, monitor appointments, review complaints, and oversee overall platform operations. \\
Doctor & View doctor-side appointments, update appointment status, write medical reports, maintain profile, and review patient complaints. \\
Patient & Register/login, manage profile, browse doctors, book appointments, cancel appointments when needed, view medical reports, and submit complaints. \\
\bottomrule
\end{tabularx}

% --------------------- 6. DETAILED MODULES ---------------------
\section{Detailed Module Description}
\subsection{Authentication Module}
\begin{itemize}[leftmargin=1.2cm]
    \item Separate registration endpoints for patient, doctor, and admin roles.
    \item Passwords are hashed before database insertion.
    \item Login checks all role tables and issues JWT with role and user id.
    \item Frontend stores session information and redirects to role dashboard.
\end{itemize}

\subsection{Profile Module}
\begin{itemize}[leftmargin=1.2cm]
    \item Authenticated users can fetch their own profile.
    \item Users can update profile information through a dedicated update endpoint.
    \item Profile pages are separated by role in the UI for clarity.
\end{itemize}

\subsection{Doctor and Patient Directory Module}
\begin{itemize}[leftmargin=1.2cm]
    \item Doctor lists are available for patient appointment booking and admin reference.
    \item Patient lists are available for admin and doctor-side monitoring.
\end{itemize}

\subsection{Appointment Module}
\begin{itemize}[leftmargin=1.2cm]
    \item Patient can create a new appointment request.
    \item Doctor can view assigned appointments and update status (for example: accepted/rejected/completed based on system flow).
    \item Patient can view personal appointments and cancel when required.
    \item Admin can view all appointments and perform administrative deletion.
\end{itemize}

\subsection{Medical Report Module}
\begin{itemize}[leftmargin=1.2cm]
    \item Doctor creates treatment/consultation reports for patients.
    \item Patient can view own medical reports from dashboard.
    \item Doctor can review reports written by self.
\end{itemize}

\subsection{Complaint Module}
\begin{itemize}[leftmargin=1.2cm]
    \item Patient can submit service complaints.
    \item Doctor can view complaints relevant to doctor-side review.
    \item Admin can view all complaints for monitoring and quality control.
\end{itemize}

\subsection{Admin Management Module}
\begin{itemize}[leftmargin=1.2cm]
    \item Consolidated user listing from role-based user tables.
    \item Endpoint support for deleting and updating user records.
    \item Role-specific edit endpoints for doctor and patient information.
\end{itemize}

% --------------------- 7. IMPLEMENTED ROUTES ---------------------
\section{Implemented API Endpoints}
\subsection{Authentication APIs}
\begin{itemize}[leftmargin=1.2cm]
    \item \texttt{POST /api/auth/register/patient}
    \item \texttt{POST /api/auth/register/doctor}
    \item \texttt{POST /api/auth/register/admin}
    \item \texttt{POST /api/auth/login}
\end{itemize}

\subsection{Profile APIs}
\begin{itemize}[leftmargin=1.2cm]
    \item \texttt{GET /api/profile/me}
    \item \texttt{PUT /api/profile/update}
\end{itemize}

\subsection{Doctor/Patient/Admin APIs}
\begin{itemize}[leftmargin=1.2cm]
    \item \texttt{GET /api/doctor/list}
    \item \texttt{GET /api/doctor/patients}
    \item \texttt{GET /api/patient/all}
    \item \texttt{GET /api/admin/users}
    \item \texttt{GET /api/admin/user/:id}
    \item \texttt{DELETE /api/admin/users/:id}
    \item \texttt{PUT /api/admin/users/:id}
    \item \texttt{PUT /api/admin/doctor/:id}
    \item \texttt{PUT /api/admin/patient/:id}
\end{itemize}

\subsection{Appointment/Report/Complaint APIs}
\begin{itemize}[leftmargin=1.2cm]
    \item \texttt{POST /api/appointment/book}
    \item \texttt{GET /api/appointment/my}
    \item \texttt{GET /api/appointment/doctor}
    \item \texttt{PUT /api/appointment/status/:id}
    \item \texttt{GET /api/appointment/all}
    \item \texttt{DELETE /api/appointment/:id}
    \item \texttt{PUT /api/appointment/cancel/:id}
    \item \texttt{POST /api/report/add}
    \item \texttt{GET /api/report/patient}
    \item \texttt{GET /api/report/doctor}
    \item \texttt{POST /api/complaint/add}
    \item \texttt{GET /api/complaint/my}
    \item \texttt{GET /api/complaint/doctor}
    \item \texttt{GET /api/complaint/all}
\end{itemize}

% --------------------- 8. DATABASE DESIGN ---------------------
\section{Database Design and Data Model}
\subsection{Database Name}
\texttt{jobra\_hospital}

\subsection{Tables in Schema}
\rowcolors{2}{SoftGray}{white}
\begin{tabularx}{\linewidth}{>{\bfseries}p{4cm} X}
\toprule
Table Name & Purpose \\
\midrule
\texttt{jh\_admins} & Stores administrator account and profile details. \\
\texttt{jh\_doctors} & Stores doctor account and professional profile details. \\
\texttt{jh\_patients} & Stores patient account and personal profile details. \\
\texttt{jh\_appointments} & Maintains appointment requests, schedule, and status. \\
\texttt{jh\_reports} & Stores clinical reports written by doctors for patients. \\
\texttt{jh\_complaints} & Stores complaint entries submitted by patients. \\
\texttt{jh\_logins} & Reserved/login-related table in schema for authentication support. \\
\bottomrule
\end{tabularx}

\subsection{Entity Relationships (Conceptual)}
\begin{itemize}[leftmargin=1.2cm]
    \item One doctor can have many appointments.
    \item One patient can have many appointments.
    \item One doctor can write many reports; one patient can receive many reports.
    \item Complaints are linked with patient and optionally doctor references.
\end{itemize}

% --------------------- 9. PROJECT WORKFLOW ---------------------
\section{End-to-End Workflow Description}
\subsection{Patient Journey}
\begin{enumerate}[leftmargin=1.2cm]
    \item Patient registers account and logs in.
    \item Patient updates profile if required.
    \item Patient checks available doctors.
    \item Patient books appointment and tracks appointment status.
    \item Patient can submit complaint and view own medical reports.
\end{enumerate}

\subsection{Doctor Journey}
\begin{enumerate}[leftmargin=1.2cm]
    \item Doctor logs in to role-specific dashboard.
    \item Doctor checks appointment requests.
    \item Doctor updates appointment status.
    \item Doctor writes medical report for consulted patient.
    \item Doctor reviews complaint and own profile details.
\end{enumerate}

\subsection{Admin Journey}
\begin{enumerate}[leftmargin=1.2cm]
    \item Admin logs in to monitoring dashboard.
    \item Admin reviews users, appointments, and complaints.
    \item Admin updates/deletes records based on management need.
    \item Admin keeps the system organized and consistent.
\end{enumerate}

% --------------------- 10. FRONTEND PAGE MAP ---------------------
\section{Frontend Route Map}
\subsection{Public Routes}
\begin{itemize}[leftmargin=1.2cm]
    \item \texttt{/} (Login)
    \item \texttt{/register}
\end{itemize}

\subsection{Patient Routes}
\begin{itemize}[leftmargin=1.2cm]
    \item \texttt{/patient}, \texttt{/patient/info}, \texttt{/patient/edit}
    \item \texttt{/patient/book}, \texttt{/patient/appointments}, \texttt{/patient/doctors}
    \item \texttt{/patient/complaint}, \texttt{/patient/report}
\end{itemize}

\subsection{Doctor Routes}
\begin{itemize}[leftmargin=1.2cm]
    \item \texttt{/doctor}, \texttt{/doctor/info}, \texttt{/doctor/edit}
    \item \texttt{/doctor/patients}, \texttt{/doctor/appointments}
    \item \texttt{/doctor/reviews}, \texttt{/doctor/write-report}
\end{itemize}

\subsection{Admin Routes}
\begin{itemize}[leftmargin=1.2cm]
    \item \texttt{/admin}, \texttt{/admin/info}, \texttt{/admin/edit}
    \item \texttt{/admin/appointments}, \texttt{/admin/add}
    \item \texttt{/admin/doctors}, \texttt{/admin/patients}, \texttt{/admin/complaints}
\end{itemize}

% --------------------- 11. INSTALLATION ---------------------
\section{Installation and Execution Guide}
\subsection{Project Structure (High Level)}
\begin{verbatim}
Jobra-Hospital-6th-Semester-Project/
├── backend/
│   ├── config/
│   ├── controllers/
│   ├── middleware/
│   ├── models/
│   ├── routes/
│   └── server.js
├── src/
│   ├── components/
│   ├── context/
│   ├── pages/
│   ├── services/
│   └── styles/
└── database.sql
\end{verbatim}

\subsection{Backend Setup}
\begin{verbatim}
cd backend
npm install
npm run dev
\end{verbatim}

\subsection{Frontend Setup}
\begin{verbatim}
npm install
npm start
\end{verbatim}

\subsection{Database Setup}
\begin{enumerate}[leftmargin=1.2cm]
    \item Create MySQL database: \texttt{jobra\_hospital}
    \item Import schema from \texttt{database.sql}
    \item Configure backend \texttt{.env} credentials
\end{enumerate}

% --------------------- 12. SECURITY NOTES ---------------------
\section{Security Features}
\begin{itemize}[leftmargin=1.2cm]
    \item Password hashing with \texttt{bcrypt}.
    \item Token-based authentication with \texttt{JWT}.
    \item Role-based authorization middleware for API access control.
    \item Frontend protected routes for role-specific navigation safety.
\end{itemize}

\subsection{Current Limitations Identified}
\begin{itemize}[leftmargin=1.2cm]
    \item JWT secret is currently hardcoded in implementation and should be moved fully to environment variables.
    \item CORS is open by default and should be restricted to approved client origins.
    \item Admin add-user payload format and register endpoint field expectations need full alignment.
\end{itemize}

% --------------------- 13. SCREENSHOT SECTION ---------------------
\section{Key Interface Screenshots}
\infobox{Replace each placeholder with your real project screenshots before final submission.}

\placeholder{Login and Register Interface}{Insert screenshot from authentication pages.}
\placeholder{Admin Dashboard}{Insert screenshot of admin home and quick actions.}
\placeholder{Doctor Dashboard}{Insert screenshot of doctor module cards/navigation.}
\placeholder{Patient Dashboard}{Insert screenshot of patient quick service panel.}
\placeholder{Book Appointment Page}{Insert screenshot of booking form and doctor selection.}
\placeholder{Doctor Appointment Management}{Insert screenshot of doctor appointment status update view.}
\placeholder{Medical Report Page}{Insert screenshot showing report writing or report list view.}
\placeholder{Complaint Module}{Insert screenshot of complaint submission/listing page.}

% --------------------- 14. TESTING SUMMARY ---------------------
\section{Testing and Validation Summary}
\subsection{Functional Validation}
\begin{itemize}[leftmargin=1.2cm]
    \item Login and role-based redirection verified.
    \item CRUD behavior validated across appointments, complaints, and reports.
    \item API request/response validation performed through frontend integration and endpoint testing.
\end{itemize}

\subsection{Usability Validation}
\begin{itemize}[leftmargin=1.2cm]
    \item Navigation is role-focused and dashboard oriented.
    \item Main workflows require limited steps and are understandable for first-time users.
\end{itemize}

% --------------------- 15. DOCUMENTATION SUITE ---------------------
\section{Documentation Suite}
\subsection{Code Documentation (GitHub Repository)}
The source code is maintained with Git and hosted in a GitHub repository for version tracking, review, and collaboration.
\begin{itemize}[leftmargin=1.2cm]
    \item Repository Link: \url{https://github.com/TanjimTajwar/Jobra-Hospital-6th-Semester-Project.git}
\end{itemize}

\subsection{Technical Documentation}
This report includes complete technical documentation, including:
\begin{itemize}[leftmargin=1.2cm]
    \item technology stack and architecture model,
    \item database schema and conceptual relationships,
    \item implemented modules and route-level API mapping,
    \item security design and system limitations.
\end{itemize}

\subsection{User Documentation}
User-facing guidance is provided through:
\begin{itemize}[leftmargin=1.2cm]
    \item installation and setup instructions,
    \item role-based workflow steps for patient, doctor, and admin users,
    \item route map and dashboard behavior explanation,
    \item screenshot placeholders for step-by-step UI demonstration.
\end{itemize}

\subsection{API Documentation}
The report contains endpoint documentation categorized by module:
\begin{itemize}[leftmargin=1.2cm]
    \item authentication, profile, admin, doctor, and patient services,
    \item appointment, report, and complaint services,
    \item access-control behavior through middleware.
\end{itemize}

% --------------------- 16. AGILE + TEAM COLLABORATION ---------------------
\section{Agile Methodology and Team Collaboration}
\subsection{Agile Methodology Adherence}
Development followed an iterative workflow inspired by agile practices:
\begin{itemize}[leftmargin=1.2cm]
    \item sprint-style module completion (auth, appointments, reports, complaints),
    \item frequent integration between frontend and backend layers,
    \item continuous fixes and refinement after feature-level testing,
    \item incremental feature addition with role-by-role delivery.
\end{itemize}

\subsection{Team Collaboration and Task Distribution}
\rowcolors{2}{SoftGray}{white}
\begin{tabularx}{\linewidth}{>{\bfseries}p{4.6cm} >{\bfseries}p{2.7cm} X}
\toprule
Member Name & ID & Contribution \\
\midrule
Tanjim Tajwar Arnab & 22701066 & Backend and database design, API development, middleware, and DB integration. \\
Sabrina Sultana & 22701067 & Admin frontend implementation, admin dashboard pages, and admin-side UI interactions. \\
Hafiz Hasnat Sifar Jami & 22701068 & Doctor frontend implementation, doctor dashboard, appointment/review/report interfaces. \\
Muznabin Ahmed & 22701069 & Patient frontend implementation, booking flow, complaint/report access, and patient UI features. \\
\bottomrule
\end{tabularx}

% --------------------- 17. ENGINEERING QUALITY EVALUATION ---------------------
\section{Engineering Quality Evaluation}
\subsection{Version Control Usage}
\begin{itemize}[leftmargin=1.2cm]
    \item Git used for source control and change history.
    \item GitHub used for repository hosting and collaboration.
    \item Commit-based progress tracking enabled safer integration.
\end{itemize}

\subsection{Feature Completeness}
\begin{itemize}[leftmargin=1.2cm]
    \item Core multi-role features are implemented end-to-end.
    \item Authentication, appointment management, reporting, complaints, and profile update are functional.
    \item Admin supervision modules are integrated for operational control.
\end{itemize}

\subsection{User Experience Quality}
\begin{itemize}[leftmargin=1.2cm]
    \item Role-focused dashboards reduce user confusion.
    \item Clear navigation paths for major workflows.
    \item Modular page design improves readability and usability.
\end{itemize}

\subsection{Error Handling}
\begin{itemize}[leftmargin=1.2cm]
    \item API-level validation and status messaging are used for major operations.
    \item Frontend handles API failures with user-visible feedback in workflow pages.
    \item Additional centralized validation and error formatting are identified as future enhancements.
\end{itemize}

\subsection{Cross-Browser Compatibility}
\begin{itemize}[leftmargin=1.2cm]
    \item Frontend uses standards-compliant React and CSS implementation.
    \item Intended compatibility: Chrome, Firefox, Edge, and other modern browsers.
    \item No browser-specific dependency is required for core workflow operation.
\end{itemize}

\subsection{Code Quality and Architecture}
\begin{itemize}[leftmargin=1.2cm]
    \item Modular structure: routes, controllers, models, services, and pages.
    \item Separation of concerns supports maintainability.
    \item Role-based middleware keeps business logic organized and secure.
\end{itemize}

\subsection{Proper Use of Chosen Technologies}
\begin{itemize}[leftmargin=1.2cm]
    \item React used for componentized role-based interfaces.
    \item Express used for scalable REST API routing.
    \item MySQL used for relational, constraint-oriented data management.
    \item JWT and bcrypt used for session control and credential security.
\end{itemize}

\subsection{Security Implementation}
\begin{itemize}[leftmargin=1.2cm]
    \item Password hashing, token authentication, and route-level role authorization are implemented.
    \item Protected frontend routing reduces unauthorized page-level access.
    \item Secret handling and stricter CORS policy are identified for hardening.
\end{itemize}

\subsection{Performance Optimization}
\begin{itemize}[leftmargin=1.2cm]
    \item Lightweight API response flow for role-specific modules.
    \item Structured database tables for direct relational querying.
    \item Scope for future optimization: pagination, indexing strategy, and caching.
\end{itemize}

% --------------------- 18. FUTURE IMPROVEMENTS ---------------------
\section{Future Scope and Enhancements}
\begin{itemize}[leftmargin=1.2cm]
    \item Notification system (email/SMS/WhatsApp) for appointment reminders.
    \item Better validation and standardized error messages for all forms.
    \item Search, filter, and pagination for large doctor/patient datasets.
    \item Prescription module and downloadable PDF reports.
    \item Advanced analytics dashboard for admin-level decision support.
    \item Stronger security hardening (Bearer token convention, stricter CORS, refresh token strategy).
\end{itemize}

% --------------------- 19. CONTRIBUTION AND REFLECTION ---------------------
\section{Contribution and Learning Reflection}
\subsection{Technical Learning}
This project strengthened understanding of:
\begin{itemize}[leftmargin=1.2cm]
    \item Full-stack architecture and module separation.
    \item Relational database design and foreign key relationships.
    \item Middleware-based authorization and secure authentication.
    \item API-driven frontend integration patterns.
\end{itemize}

\subsection{Project Outcome Reflection}
The completed system serves as a practical and scalable foundation for a real hospital management product. It demonstrates that structured DBMS design plus modular full-stack development can significantly improve service quality, record consistency, and administrative control.

% --------------------- 20. CONCLUSION ---------------------
\section{Conclusion}
The \textbf{Jobra Hospital Management System} successfully fulfills its primary mission: transforming manual hospital operations into a digital, role-based, data-centric workflow. From login and profile management to appointments, reporting, complaints, and centralized administration, the application provides an academically strong and functionally complete demonstration of DBMS-centered full-stack engineering.

The current implementation is deployment-ready for academic demonstration and can be elevated to production quality through incremental improvements in security hardening, validation quality, and operational tooling.

\section*{License}
\addcontentsline{toc}{section}{License}
This software is developed as an academic project for educational purposes.

\end{document}